\documentclass[conference]{IEEEtran}
\IEEEoverridecommandlockouts

\usepackage{cite}
\usepackage{amsmath,amssymb,amsfonts}
\usepackage{algorithmic}
\usepackage{graphicx}
\usepackage{textcomp}
\usepackage{xcolor}
\usepackage{listings}
\usepackage{booktabs}
\usepackage{url}

\lstset{
    basicstyle=\footnotesize\ttfamily,
    frame=single,
    breaklines=true,
    numbers=left,
    numberstyle=\tiny\color{gray},
    keywordstyle=\color{blue}\bfseries,
    commentstyle=\color{gray}\itshape,
    stringstyle=\color{purple}
}

\def\BibTeX{{\rm B\kern-.05em{\sc i\kern-.025em b}\kern-.08em
    T\kern-.1667em\lower.7ex\hbox{E}\kern-.125emX}}

\begin{document}

\title{Digital Cookbook: An Interactive Domain-Specific Language Compiler for Culinary Automation and Static Verification}

\author{
\IEEEauthorblockN{[Student 1 Name - Team Lead]}
\IEEEauthorblockA{\textit{Dept. of Computer Science \& Engineering} \\
\textit{[University / Institute Name]}\\
\textit{Roll No: [Roll Number 1]} \\
email@institution.edu}
\and
\IEEEauthorblockN{[Student 2 Name]}
\IEEEauthorblockA{\textit{Dept. of Computer Science \& Engineering} \\
\textit{[University / Institute Name]}\\
\textit{Roll No: [Roll Number 2]} \\
email@institution.edu}
\and
\IEEEauthorblockN{[Student 3 Name]}
\IEEEauthorblockA{\textit{Dept. of Computer Science \& Engineering} \\
\textit{[University / Institute Name]}\\
\textit{Roll No: [Roll Number 3]} \\
email@institution.edu}
\and
\IEEEauthorblockN{[Student 4 Name]}
\IEEEauthorblockA{\textit{Dept. of Computer Science \& Engineering} \\
\textit{[University / Institute Name]}\\
\textit{Roll No: [Roll Number 4]} \\
email@institution.edu}
}

\maketitle

\begin{abstract}
Traditional natural-language culinary recipes are ambiguous, non-standardized, and vulnerable to measurement, sequencing, and semantic contradictions. In modern commercial kitchens and autonomous robotic culinary systems, precise verification and deterministic machine instruction generation are essential. This paper presents \textbf{Digital Cookbook}, an interactive, web-based compiler for a cooking-oriented domain-specific language termed \textbf{Cookbook DSL}. The system implements a complete, decoupled six-phase compiler pipeline: (1) a hand-written character-level scanner with precise coordinate tracking; (2) a recursive-descent syntax parser constructing a strongly-typed Abstract Syntax Tree (AST); (3) a domain-aware semantic analyzer and symbol table enforcing recipe consistency, ingredient compatibility, and execution order; (4) a Three-Address Code (3AC) intermediate representation generator; (5) a multi-pass optimizer performing measurement normalization, adjacent instruction aggregation, and dead-code elimination; and (6) a target code generator emitting Virtual Kitchen Machine (VKM) bytecode assembly. A comprehensive evaluation comprising 42 automated test cases proves the deterministic behavior, rigorous diagnostic reporting, and performance efficiency of the architecture.
\end{abstract}

\begin{IEEEkeywords}
Compiler Design, Domain-Specific Languages, Lexical Analysis, Recursive-Descent Parser, Abstract Syntax Tree, Semantic Analysis, Intermediate Representation, Peephole Optimization, Code Generation.
\end{IEEEkeywords}

\section{Introduction}
Compilers serve as the fundamental backbone of computer science, translating human-intelligible linguistic constructs into deterministic hardware execution models. While general-purpose languages such as C, Java, and Python prioritize general arithmetic and memory operations, Domain-Specific Languages (DSLs) provide tailored expressive power for specialized application domains. 

In automated culinary manufacturing and intelligent food robotics, natural-language recipes introduce significant hazards: vague volumetric units, unspecified dependencies, missing safety steps, and undefined heating profiles. To address these challenges, we introduce \textbf{Digital Cookbook}, an end-to-end interactive compiler engineering project designed to compile, verify, optimize, and assemble structured cooking instructions into low-level Virtual Kitchen Machine instructions.

\section{Problem Statement}
Existing recipe representation formats suffer from three critical defects:
\begin{enumerate}
    \item \textbf{Ambiguity and Non-determinism}: Vague directives such as ``cook until ready'' prevent automated robotic reproduction.
    \item \textbf{Absence of Static Verification}: Semantic errors—such as baking an empty vessel, combining incompatible ingredients, or exceeding physical appliance thresholds—are detected only at runtime or after ingredient spoilage.
    \item \textbf{Inefficient Actuation}: Sequential recipe declarations often describe fragmented ingredient additions that trigger redundant mechanical actuation cycles unless aggregated via formal compiler optimization passes.
\end{enumerate}

\section{System Architecture}
The compiler adheres to a classic decoupled six-phase architecture, orchestrated by a central pipeline service. The system is entirely self-contained, deterministic, and free of external AI dependencies.

\begin{figure}[htbp]
\centering
\fbox{\parbox{0.9\columnwidth}{
\centering
\textbf{Cookbook DSL Source Code}\\
$\downarrow$\\
\textbf{Phase 1: Lexical Analysis} (Character Scanner $\rightarrow$ Token Stream)\\
$\downarrow$\\
\textbf{Phase 2: Syntax Analysis} (Recursive Descent $\rightarrow$ Typed AST)\\
$\downarrow$\\
\textbf{Phase 3: Semantic Analysis} (Context Validation $\rightarrow$ Symbol Table)\\
$\downarrow$\\
\textbf{Phase 4: Intermediate Code Generation} (Linear Three-Address Code)\\
$\downarrow$\\
\textbf{Phase 5: Code Optimization} (Peephole Merging \& Unit Normalization)\\
$\downarrow$\\
\textbf{Phase 6: Target Code Generation} (Virtual Kitchen Machine Assembly)
}}
\caption{Digital Cookbook Compiler Pipeline Architecture.}
\label{fig:architecture}
\end{figure}

\section{Cookbook DSL Language Specification}
Cookbook DSL is a declarative, procedural domain-specific language. Every program models one complete culinary preparation.

\subsection{Formal EBNF Grammar}
The formal syntax is specified in Extended Backus-Naur Form (EBNF):
\begin{align*}
\text{program} &\to \text{recipe\_decl}\;\;\{\text{statement}\}\;\;\text{serve\_stmt}\;\;\text{EOF} \\
\text{recipe\_decl} &\to \text{"RECIPE"}\;\;\text{identifier}\;\;\text{";"} \\
\text{statement} &\to \text{add\_stmt} \mid \text{timed\_stmt} \mid \text{chop\_stmt} \\
\text{add\_stmt} &\to \text{"ADD("}\;\text{number}\text{","}\;\text{unit}\text{","}\;\text{identifier}\text{");"} \\
\text{timed\_stmt} &\to \text{timed\_op}\text{"("}\;\text{number}\text{","}\;\text{time\_unit}\text{");"} \\
\text{serve\_stmt} &\to \text{"SERVE;"}
\end{align*}

\section{Compiler Implementation Phases}

\subsection{Phase 1: Lexical Analysis}
The lexer (\texttt{compiler/lexer.py}) scans UTF-8 source text character-by-character. It preserves token lexemes while accurately tracking line and column numbers. Reserved keywords (\texttt{RECIPE}, \texttt{ADD}, \texttt{MIX}, \texttt{BOIL}, \texttt{BAKE}, \texttt{FRY}, \texttt{SERVE}), measurement units (\texttt{CUP}, \texttt{GRAM}, \texttt{ML}, \texttt{PIECE}), and time units (\texttt{SECONDS}, \texttt{MINUTES}) are categorized. Illegal characters are captured as \texttt{UNKNOWN} tokens and reported via structured diagnostics without crashing.

\subsection{Phase 2: Syntax Analysis}
The parser (\texttt{compiler/parser.py}) applies top-down recursive descent with panic-mode error recovery. It enforces strict statement structures and constructs an Abstract Syntax Tree composed of strongly-typed dataclasses (\texttt{ProgramNode}, \texttt{AddNode}, \texttt{TimedOpNode}, \texttt{ServeNode}). When a syntax error occurs, the parser logs the diagnostic and synchronizes to the next statement delimiter.

\subsection{Phase 3: Semantic Analysis \& Symbol Table}
The semantic analyzer (\texttt{compiler/semantic.py}) inspects the AST against a centralized culinary knowledge base. It enforces twelve domain rules:
\begin{enumerate}
    \item Verification of dish declaration against supported recipes.
    \item Ingredient validation and compatibility with the declared dish.
    \item Physical unit compatibility (e.g., eggs must be in pieces, liquids in milliliters).
    \item Strictly positive quantity and duration constraints ($q > 0, t > 0$).
    \item Operational bounds on appliance durations.
    \item Prohibition of empty-vessel operations.
    \item Mandatory thermal preparation prior to serving cooked dishes.
    \item Enforcement of terminal \texttt{SERVE} positioning.
\end{enumerate}
The symbol table maintains comprehensive records tracking identifier categories, cumulative quantities, and source locations.

\subsection{Phase 4: Intermediate Code Generation (IR)}
The intermediate representation generator (\texttt{compiler/ir.py}) linearizes the validated AST into machine-independent Three-Address Code (3AC):
\begin{lstlisting}[language=Python]
t1 = RECIPE PANCAKE
t2 = ADD FLOUR, 150.0 GRAM
t3 = ADD FLOUR, 150.0 GRAM
t4 = ADD MILK, 1.0 CUP
t5 = MIX 2.0 MINUTES
t6 = SERVE
\end{lstlisting}

\subsection{Phase 5: Code Optimization}
The optimizer (\texttt{compiler/optimizer.py}) performs three semantic-preserving transformation passes:
\begin{enumerate}
    \item \textbf{Unit Normalization}: Standardizes volume units (\texttt{CUP}, \texttt{TBSP}, \texttt{TSP}) into base milliliters (\texttt{ML}), mass into \texttt{GRAM}, and duration into \texttt{SECONDS}.
    \item \textbf{Adjacent Instruction Merging}: Combines contiguous additions of identical ingredients without intermediate processing into a single instruction. For instance, consecutive $150\,\text{g}$ additions of flour become a single $300\,\text{g}$ addition.
    \item \textbf{Dead Code Elimination}: Prunes null-effect instructions.
\end{enumerate}

\subsection{Phase 6: Target Code Generation (VKM Assembly)}
The code generator translates optimized IR into bytecode for the Virtual Kitchen Machine (VKM):
\begin{lstlisting}
    0x0000    LOAD_VESSEL PANCAKE
    0x0001    DISPENSE FLOUR, 300.0, GRAM
    0x0002    DISPENSE MILK, 240.0, ML
    0x0003    STIR 120, 120s
    0x0004    PLATING_SERVE DEFAULT
    0x0005    HALT
\end{lstlisting}

\section{Testing and Verification}
Quality assurance was conducted using \texttt{pytest}. A test suite comprising 42 distinct unit and integration test cases was executed across all compiler phases:
\begin{itemize}
    \item Lexer coordinate tracking and malformed numeral detection.
    \item Parser grammar validation and recovery from missing delimiters.
    \item Semantic multi-error reporting and symbol table accumulation.
    \item Intermediate representation consistency.
    \item Optimization semantic equivalence and instruction reduction.
    \item Target assembly cycle estimation.
    \item Flask REST API endpoint verification.
\end{itemize}
All 42 test cases passed with 100\% success rate in 0.08 seconds.

\section{Individual Team Contributions}
\begin{itemize}
    \item \textbf{[Student 1 Name - Team Lead]}: Designed compiler architecture, lexical analyzer, token structures, and central pipeline service.
    \item \textbf{[Student 2 Name]}: Defined formal EBNF grammar, recursive-descent parser, AST node hierarchy, and syntax error synchronization.
    \item \textbf{[Student 3 Name]}: Implemented semantic validation rules, domain knowledge base, and symbol table manager.
    \item \textbf{[Student 4 Name]}: Developed Three-Address Code IR generation, multi-pass optimization engine, and VKM code generator.
\end{itemize}

\section{Conclusion and Future Scope}
The Digital Cookbook project successfully demonstrates the application of formal compiler engineering techniques to domain-specific culinary automation. Future extensions include concurrent multi-vessel scheduling, physical temperature thermodynamic modeling, and hardware synthesis for real-world automated kitchen actuators.

\section*{References}
\begin{thebibliography}{00}
\bibitem{aho} A. V. Aho, M. S. Lam, R. Sethi, and J. D. Ullman, \textit{Compilers: Principles, Techniques, and Tools}, 2nd ed. Boston, MA: Addison-Wesley, 2006.
\bibitem{cooper} K. D. Cooper and L. Torczon, \textit{Engineering a Compiler}, 2nd ed. Burlington, MA: Morgan Kaufmann, 2011.
\bibitem{muchnick} S. S. Muchnick, \textit{Advanced Compiler Design and Implementation}. San Francisco, CA: Morgan Kaufmann, 1997.
\bibitem{reference_repo} K. Yadav, ``Cookbook Compiler: A Cooking-Oriented DSL Implementation,'' GitHub Repository, 2024. [Online]. Available: \url{https://github.com/antedotee/cookbook-compiler}.
\end{thebibliography}

\end{document}
